% GENERATED FILE. Edit canonical lessons, then run npm run generate:books.
% canonical-source-hash: fnv1a64:0643906f7b87a713
% canonical-lessons: PT-C465-rajada, PT-C465-soprar, PT-C465-temporal, PT-C465-furacao, PT-C465-nevasca

\chapter{Vento e tempestade}
\label{ch:pt-vento-e-tempestade}
\hlchaptermodality{465}

\begin{chapteropening}
\textbf{By the end of this chapter:} \emph{I can describe wild weather: a gust of wind, the wind blowing, a storm, a hurricane and a blizzard.}

Complete the last lesson of chapter 465: I can describe wild weather: a gust of wind, the wind blowing, a storm, a hurricane and a blizzard.
\end{chapteropening}

\section[a rajada]{a rajada — the gust (of wind)}

\label{lesson:PT-C465-rajada}



\begin{quote}
Before the new one: say the Portuguese for the freezer, then the Portuguese for to melt.
\end{quote}

\subsection*{You'll want to know: a rajada}
\textbf{a rajada} — \textquotedblleft{}the gust (of wind)\textquotedblright{}.

\begin{quote}
Uma rajada de vento levou o meu chapéu.
\end{quote}

\begin{quote}
\emph{A gust of wind took my hat.}
\end{quote}

Forecasts give \emph{rajadas} in kilometres per hour.

\begin{cousinweb}
The origin of \emph{rajada} is not settled, and EtymDB has no entry for it.
\end{cousinweb}

\begin{grammarlens}[title={What you've built}]
One more word for telling, explaining and describing at B1.
\end{grammarlens}

\subsection*{Your turn}
\begin{itemize}
\raggedright
  \item \emph{Say it:} \emph{a rajada}
  \item \emph{Say it:} \emph{a rajada}, once more
  \item \emph{Say it:} a sentence of your own with \emph{a rajada}
  \item \emph{Recall:} say the Portuguese for the freezer, then the Portuguese for to melt, then say \emph{a rajada} again
\end{itemize}

\begin{tcolorbox}[breakable,colback=teal!4,colframe=teal!35!black,title={Before you move on}]
What does a rajada mean? (\textquotedblleft{}The gust (of wind)\textquotedblright{}.) Say it once more.
\end{tcolorbox}
\section[soprar]{soprar — to blow}

\label{lesson:PT-C465-soprar}



\begin{quote}
Before the new one: say the Portuguese for to melt, then the Portuguese for the gust.
\end{quote}

\subsection*{You'll want to know: soprar}
\textbf{soprar} — \textquotedblleft{}to blow\textquotedblright{}.

\begin{quote}
O vento sopra forte do mar.
\end{quote}

\begin{quote}
\emph{The wind is blowing hard off the sea.}
\end{quote}

People \emph{sopram} the candles on a cake too.

\begin{cousinweb}
\emph{soprar} is inherited from Latin \emph{sufflāre} 'to blow on' (EtymDB), built on \emph{flāre} 'to blow'.
\end{cousinweb}

\begin{grammarlens}[title={What you've built}]
One more word for telling, explaining and describing at B1.
\end{grammarlens}

\subsection*{Your turn}
\begin{itemize}
\raggedright
  \item \emph{Say it:} \emph{soprar}
  \item \emph{Say it:} \emph{soprar}, once more
  \item \emph{Say it:} a sentence of your own with \emph{soprar}
  \item \emph{Recall:} say the Portuguese for to melt, then the Portuguese for the gust, then say \emph{soprar} again
\end{itemize}

\begin{tcolorbox}[breakable,colback=teal!4,colframe=teal!35!black,title={Before you move on}]
What does soprar mean? (\textquotedblleft{}To blow\textquotedblright{}.) Say it once more.
\end{tcolorbox}
\section[o temporal]{o temporal — the storm (wind and rain)}

\label{lesson:PT-C465-temporal}



\begin{quote}
Before the new one: say the Portuguese for the gust, then the Portuguese for to blow.
\end{quote}

\subsection*{You'll want to know: o temporal}
\textbf{o temporal} — \textquotedblleft{}the storm (wind and rain)\textquotedblright{}.

\begin{quote}
O temporal deixou muitas árvores no chão.
\end{quote}

\begin{quote}
\emph{The storm left many trees on the ground.}
\end{quote}

A long spell of bad weather, often in winter.

\begin{cousinweb}
\emph{temporal} 'storm' is built on \emph{tempo} 'weather' (EtymDB), which is inherited from Latin \emph{tempus} 'time, season'.
\end{cousinweb}

\begin{grammarlens}[title={What you've built}]
One more word for telling, explaining and describing at B1.
\end{grammarlens}

\subsection*{Your turn}
\begin{itemize}
\raggedright
  \item \emph{Say it:} \emph{o temporal}
  \item \emph{Say it:} \emph{o temporal}, once more
  \item \emph{Say it:} a sentence of your own with \emph{o temporal}
  \item \emph{Recall:} say the Portuguese for the gust, then the Portuguese for to blow, then say \emph{o temporal} again
\end{itemize}

\begin{tcolorbox}[breakable,colback=teal!4,colframe=teal!35!black,title={Before you move on}]
What does o temporal mean? (\textquotedblleft{}The storm (wind and rain)\textquotedblright{}.) Say it once more.
\end{tcolorbox}
\section[o furacão]{o furacão — the hurricane}

\label{lesson:PT-C465-furacao}



\begin{quote}
Before the new one: say the Portuguese for to blow, then the Portuguese for the storm.
\end{quote}

\subsection*{You'll want to know: o furacão}
\textbf{o furacão} — \textquotedblleft{}the hurricane\textquotedblright{}.

\begin{quote}
O furacão passou longe da costa.
\end{quote}

\begin{quote}
\emph{The hurricane passed far from the coast.}
\end{quote}

Figurative: a very lively person.

\begin{cousinweb}
\emph{furacão} was borrowed from Spanish \emph{huracán}, which comes from Taíno \emph{hurakán}, a word of the Caribbean islands (EtymDB).
\end{cousinweb}

\begin{grammarlens}[title={What you've built}]
One more word for telling, explaining and describing at B1.
\end{grammarlens}

\subsection*{Your turn}
\begin{itemize}
\raggedright
  \item \emph{Say it:} \emph{o furacão}
  \item \emph{Say it:} \emph{o furacão}, once more
  \item \emph{Say it:} a sentence of your own with \emph{o furacão}
  \item \emph{Recall:} say the Portuguese for to blow, then the Portuguese for the storm, then say \emph{o furacão} again
\end{itemize}

\begin{tcolorbox}[breakable,colback=teal!4,colframe=teal!35!black,title={Before you move on}]
What does o furacão mean? (\textquotedblleft{}The hurricane\textquotedblright{}.) Say it once more.
\end{tcolorbox}
\section[a nevasca]{a nevasca — the blizzard, the snowstorm}

\label{lesson:PT-C465-nevasca}



\begin{quote}
Before the new one: say the Portuguese for the storm, then the Portuguese for the hurricane.
\end{quote}

\subsection*{You'll want to know: a nevasca}
\textbf{a nevasca} — \textquotedblleft{}the blizzard, the snowstorm\textquotedblright{}.

\begin{quote}
A nevasca fechou a estrada da serra.
\end{quote}

\begin{quote}
\emph{The blizzard closed the mountain road.}
\end{quote}

Built on \emph{a neve}.

\begin{cousinweb}
\emph{nevasca} is built on \emph{neve} 'snow', which is inherited from Latin \emph{nix, nivis} (EtymDB).
\end{cousinweb}

\begin{grammarlens}[title={What you've built}]
One more word for telling, explaining and describing at B1.
\end{grammarlens}

\subsection*{Your turn}
\begin{itemize}
\raggedright
  \item \emph{Say it:} \emph{a nevasca}
  \item \emph{Say it:} \emph{a nevasca}, once more
  \item \emph{Say it:} a sentence of your own with \emph{a nevasca}
  \item \emph{Recall:} say the Portuguese for the storm, then the Portuguese for the hurricane, then say \emph{a nevasca} again
\end{itemize}

\begin{tcolorbox}[breakable,colback=teal!4,colframe=teal!35!black,title={Before you move on}]
What does a nevasca mean? (\textquotedblleft{}The blizzard, the snowstorm\textquotedblright{}.) Say it once more.
\end{tcolorbox}
